\documentclass[11pt,a4paper,sans]{moderncv}

\moderncvstyle{banking}
\moderncvcolor{blue}

\usepackage[utf8]{inputenc}
\usepackage[T1]{fontenc}
\usepackage{lmodern}
\usepackage[scale=0.85]{geometry}
\renewcommand*{\namefont}{\fontsize{28}{32}\mdseries\upshape}
\renewcommand*{\titlefont}{\fontsize{22}{26}\mdseries\upshape}

\name{\.{I}brahim}{Sayar}
\title{K{\i}demli Yaz{\i}l{\i}m M\"uhendisi}
\address{Augsburg, Almanya}{}{}
\phone[mobile]{+49~170~1030979}
\email{hello@ibrahimsayar.com}
\homepage{ibrahimsayar.com}
\social[linkedin]{ibrahimsayar}
\social[github]{ibrahimsayar}

\begin{document}

\makecvtitle

\vspace{-2em}

%-----------------------------------------------------------------------
% HAKKIMDA
%-----------------------------------------------------------------------
\section{Hakk{\i}mda}
Yaz{\i}l{\i}m mimarisi olu\c{s}turma, backend geli\c{s}tirme ve DevOps alanlar{\i}nda 10 y{\i}l{\i} a\c{s}k{\i}n deneyime sahip k{\i}demli bir yaz{\i}l{\i}m m\"uhendisiyim. Dahil oldu\u{g}um sistemlerin i\c{s}leyi\c{s}ini h{\i}zla kavrama ve projelere k{\i}sa s\"urede uyum sa\u{g}lama yetene\u{g}ine sahibim. Uzmanl{\i}k alanlar{\i}mdan biri, kar\c{s}{\i}la\c{s}t{\i}\u{g}{\i}n{\i}z sorunlara en iyi uygulama (best practice) standartlar{\i}na uygun teknik \"c\"oz\"umlerle yan{\i}t verebilmektir. Kariyerim boyunca y\"uksek trafikli mimarilerde g\"orev ald{\i}m; bu s\"ure\c{c}te sorunlar{\i}n \"c\"oz\"um\"un\"u, sistem iyile\c{s}tirmelerini ve yeni \"ozelliklerin geli\c{s}tirilmesini u\c{c}tan uca \"ustlenerek ba\c{s}ar{\i}l{\i} sonu\c{c}lar ortaya koydum. Kendimi olduk\c{c}a yetkin g\"or\"uyor; sorumluluk almay{\i} ve y\"uksek performans sergilemeyi gerektiren, tempolu \c{c}al{\i}\c{s}ma ortamlar{\i}nda ba\c{s}ar{\i}l{\i} bir \c{s}ekilde \c{c}al{\i}\c{s}abiliyorum. Karma\c{s}{\i}k ve da\u{g}{\i}n{\i}k kod tabanlar{\i}n{\i} okumaktan, i\c{c}inde gezinmekten ve d\"uzenlemekten \c{c}ekinmem --- \"uretken olmak i\c{c}in temiz koda ihtiyac{\i}m yok.

%-----------------------------------------------------------------------
% DENEYIM
%-----------------------------------------------------------------------
\section{Deneyim}

\cventry{May{\i}s 2026 -- Halen}{K{\i}demli Yaz{\i}l{\i}m M\"uhendisi}{Naylalabs}{Uzaktan}{}{
\begin{itemize}
\item Ayl{\i}k 750.000'den fazla sipari\c{s} i\c{s}leyen \c{c}e\c{s}itli e-ticaret altyap{\i}lar{\i}n{\i}n (API, web ve pazar yeri) backend mimarisi ve DevOps s\"ure\c{c}lerinden sorumluyum.
\item E-ticaret platformlar{\i}, ERP yaz{\i}l{\i}mlar{\i} ve kurumsal y\"onetim sistemleri dahil 100'den fazla sistemin mimari y\"onlerini belirliyorum. ERP sistemlerini ve hem dahili hem harici servisleri, domain-driven design s{\i}n{\i}rlar{\i}na s{\i}k{\i} s{\i}k{\i}ya ba\u{g}l{\i} kalarak olay g\"ud\"uml\"u entegrasyonlarla birbirine ba\u{g}l{\i}yorum.
\item Sunucular, y\"uk dengeleyiciler, CDN'ler, veritaban{\i} shard'lar{\i} ve izleme altyap{\i}s{\i}n{\i} kapsayan \"uretim ortam{\i}n{\i} y\"onetiyor ve izliyorum. Yaln{\i}zca bu e-ticaret altyap{\i}lar{\i}n{\i}n en b\"uy\"u\u{g}\"u i\c{c}in ayl{\i}k 1 milyar API iste\u{g}i i\c{s}liyorum.
\end{itemize}
}

\cventry{Ocak 2025 -- May{\i}s 2026}{TestOps M\"uhendisi}{Insider One}{T\"urkiye}{}{
\begin{itemize}
\item Haftal{\i}k 7 milyar{\i}n \"uzerinde istek i\c{s}leyen mikroservisler i\c{c}in, bir Kubernetes ekosistemi dahilinde CI/CD s\"ure\c{c}lerinden, test ortamlar{\i}ndan ve bu test ortamlar{\i}n{\i} destekleyen altyap{\i}dan (sunuculardan) sorumluydum.
\item Yaz{\i}l{\i}m kalitesi, performans ve iyile\c{s}tirmeler odakl{\i} ek destek sa\u{g}lad{\i}m. Ayr{\i}ca ger\c{c}ek zamanl{\i} izleme ara\c{c}lar{\i} kullanarak analizler ger\c{c}ekle\c{s}tirdim ve bulgular{\i} ilgili ekiplere ilettim.
\item AWS Cloud Console altyap{\i}s{\i} dahilindeki t\"um \"ur\"un ve servisleri y\"onettim ve yap{\i}land{\i}rd{\i}m.
\end{itemize}
}

\cventry{A\u{g}ustos 2023 -- \c{S}ubat 2025}{DevOps M\"uhendisi}{PangeaVR}{\.{I}stanbul, T\"urkiye}{}{
\begin{itemize}
\item AWS Cloud Console sistemlerini kullanarak, Kubernetes altyap{\i}s{\i} ile yo\u{g}un oyun isteklerini kar\c{s}{\i}l{\i}yor, monit\"or ediyor ve bunlar{\i}n kurulum, bak{\i}m, geli\c{s}tirmelerini sa\u{g}l{\i}yordum.
\item Development'tan production a\c{s}amas{\i}na kadar, CI/CD s\"ure\c{c}lerini y\"onetiyor ve aksamamalar{\i}n{\i} sa\u{g}l{\i}yordum.
\end{itemize}
}

\cventry{A\u{g}ustos 2022 -- A\u{g}ustos 2023}{Full-stack Geli\c{s}tirici \& DevOps}{Naylalabs}{Samsun, T\"urkiye}{}{
Node.js, Python ve JavaScript ile full-stack web uygulamalar{\i} geli\c{s}tirme.
RESTful API tasar{\i}m{\i} ve PostgreSQL/MySQL veritaban{\i} y\"onetimi.
Docker ve CI/CD otomasyonu ile konteynerize da\u{g}{\i}t{\i}m uygulamalar{\i}.
}

\cventry{Temmuz 2021 -- Eyl\"ul 2022}{Web Geli\c{s}tirici}{Naylalabs}{Samsun, T\"urkiye}{}{
Modern JavaScript framework'leri ile duyarl{\i} web uygulamalar{\i} geli\c{s}tirme.
E-posta servisleri ve JSON API'leri dahil \"u\c{c}\"unc\"u parti hizmet entegrasyonlar{\i}.
Git s\"ur\"um kontrol\"u ile \c{c}evik i\c{s} ak{\i}\c{s}lar{\i}nda i\c{s} birli\u{g}i.
}

\cventry{Ocak 2020 -- Haziran 2021}{K{\i}demli Back-end Geli\c{s}tirici}{Veriweb}{Samsun, T\"urkiye}{}{
Back-end geli\c{s}tirme ekibinin y\"onetimi, yaz{\i}l{\i}m mimarisi ve kodlama standartlar{\i}n{\i}n belirlenmesi.
PHP ve Python ile \"ol\c{c}eklenebilir sunucu taraf{\i} uygulamalar{\i} geli\c{s}tirme.
SQL veritabanlar{\i}n{\i}n (MySQL, PostgreSQL) tasarlanmas{\i} ve optimizasyonu.
}

\cventry{Kas{\i}m 2016 -- Aral{\i}k 2019}{Kurucu Ortak \& Back-end Geli\c{s}tirici}{Kodakil}{Samsun, T\"urkiye}{}{
Bir teknoloji giri\c{s}iminin kurucu ortakl{\i}\u{g}{\i}, back-end mimari ve geli\c{s}tirme liderli\u{g}i.
RESTful API geli\c{s}tirme ve Linux \"uzerinde sunucu altyap{\i}s{\i} y\"onetimi.
Teknik strateji ve \"ur\"un geli\c{s}tirme ya\c{s}am d\"ong\"us\"un\"un y\"onetimi.
}

\cventry{May{\i}s 2014 -- Kas{\i}m 2016}{Web Tasar{\i}m, Programlama \& BT}{Fikir Kul\"ub\"u}{Samsun, T\"urkiye}{}{
PHP ve JavaScript ile web uygulamalar{\i} geli\c{s}tirme ve bak{\i}m{\i}.
BT altyap{\i}s{\i} y\"onetimi ve teknik destek.
}

\cventry{\c{S}ubat 2013 -- \c{S}ubat 2014}{Web Geli\c{s}tirici}{Teknocenter}{}{}{
HTML, CSS, JavaScript ve PHP ile web sitesi geli\c{s}tirme.
}

%-----------------------------------------------------------------------
% TEKNIK BECERILER
%-----------------------------------------------------------------------
\section{Teknik Beceriler}

\cvitem{Mimari}{Sistem mimarisi, Domain-Driven Design, mikroservisler, olay g\"ud\"uml\"u \& kuyruk tabanl{\i} sistemler, sistem entegrasyonu (pazar yerleri, ERP), eski sistem modernizasyonu, \"onbellekleme, y\"uk dengeleme}
\cvitem{Diller}{PHP, TypeScript, JavaScript, Python, SQL, Bash}
\cvitem{Back-end}{Laravel, Node.js, NestJS, FastAPI, REST API'leri, BullMQ, Redis}
\cvitem{Front-end}{React, Next.js, Vite, HTML, CSS}
\cvitem{DevOps \& Bulut}{GitHub Actions, Docker, Kubernetes, Terraform, AWS, Linux, nginx, Cloudflare, CI/CD}
\cvitem{\.{I}zleme \& G\"uvenlik}{Grafana, OAuth 2.0/JWT, Cloudflare WAF, SSL/TLS}
\cvitem{Veritabanlar{\i}}{PostgreSQL, MySQL, Redis, MongoDB, Oracle (entegrasyon)}
\cvitem{Test}{Selenium, Appium, Postman, u\c{c}tan uca otomasyon}

%-----------------------------------------------------------------------
% EGITIM
%-----------------------------------------------------------------------
\section{E\u{g}itim}

\cventry{2017 -- 2023}{Lisans -- Y\"onetim Bili\c{s}im Sistemleri}{Anadolu \"Universitesi}{}{}{Denklik de\u{g}erlendirmesi (ZAB): ba\c{s}vurusu yap{\i}ld{\i}}
\cventry{2014 -- 2016}{\"On Lisans -- Bilgisayar Programc{\i}l{\i}\u{g}{\i}}{Ondokuz May{\i}s \"Universitesi}{}{}{}
\cventry{2011 -- 2014}{Bilgisayar Programc{\i}l{\i}\u{g}{\i}}{Atakum Teknik ve End\"ustri Meslek Lisesi}{}{}{}

%-----------------------------------------------------------------------
% DILLER
%-----------------------------------------------------------------------
\section{Diller}

\cvitemwithcomment{T\"urk\c{c}e}{Anadil}{}
\cvitemwithcomment{\.{I}ngilizce}{B1 (Orta)}{}
\cvitemwithcomment{Almanca}{A2 (Entegrasyon kursu -- Integrationskurs, devam ediyor)}{}

%-----------------------------------------------------------------------
% ODULLLER
%-----------------------------------------------------------------------
\section{\"Od\"uller}

\cvitem{}{Ondokuz May{\i}s \"Universitesi -- Y\"uksek Onur Belgesi}

\end{document}
